\section{Physical Analysis}\label{sec:physics}

This section examines how body morphology changes the physical control
demand of a motion, using the simulator's instrumented quantities and not
claims about human biomechanics. The question is not whether body mass or
size correlates with joint torque, which is partly expected from mechanics,
but whether the effect of morphology on control demand \emph{depends on the
motion being performed}.

The procedure holds the motion fixed and varies the body. For the full
matched corpus of 150 clips evaluated across all 128 trained shape
configurations (19,200 deterministic rollouts total, epoch-28,170
checkpoint, zero failed episodes), we record per-joint torque, joint power,
an aggregate control-effort measure, and jerk from the simulator. Clips are
grouped by their inherited text descriptions into coarse semantic classes
(walking, running, jumping, kicking, dancing, sitting, standing, plus an
unclassified \emph{other} bucket) via keyword matching on the concatenated
captions, and we compare the cross-morphology variation of each quantity
between classes.

The hypothesis under test is that morphological variation produces
motion-dependent changes in control demand, with highly dynamic motions
(running, jumping, kicking) more sensitive to body shape than slow or
quasi-static motions (standing, sitting) -- that is, body morphology and
motion characteristics interact in setting control demand, rather than
morphology acting as a motion-independent scaling.

\paragraph{Instrumentation validation.} A first pass compared each class's
per-shape coefficient of variation (CV, std/mean). This produced a
highly significant difference ($p<10^{-15}$ by Levene's test on every
metric) -- but in the direction \emph{opposite} the hypothesis, with
quasi-static classes appearing more shape-sensitive than dynamic ones. The
cause is a normalization artifact rather than a real effect: quasi-static
clips have a much smaller mean control demand, so dividing by that small
mean inflates relative variance from ordinary simulation noise. We therefore
report absolute, unit-preserving residuals instead (each clip's value minus
its own across-shape mean, not divided by it), and we perform inference at
the clip level rather than the row level: a clip's 128 shape replicates are
correlated observations of that one clip, not 128 independent draws about
its semantic class, so a naive per-row test overstates significance
(pseudo-replication).

\begin{table}[t]
\centering
\small
\caption{Cross-shape variability of control demand, dynamic
(running~+~jumping~+~kicking, $n{=}25$ clips) vs.\ quasi-static
(standing~+~sitting, $n{=}37$ clips) motion classes. Each cell is the
standard deviation of the pooled, per-clip-demeaned residual (absolute
units, clip baseline removed) with a 95\% bootstrap CI resampled at the
clip level (10,000 resamples). $p$ is Levene's test (Brown--Forsythe,
median-centered) on the same pooled residuals.}\label{tab:physical_analysis}
\begin{tabular}{lrrr}
\toprule
Metric & Dynamic std [95\% CI] & Quasi-static std [95\% CI] & Levene's $p$ \\
\midrule
Normalized torque (frac.\ effort limit) & $1.93$ [$1.71$, $2.21$] $\times10^{-4}$ & $2.06$ [$1.86$, $2.30$] $\times10^{-4}$ & $0.10$ \\
Mass-normalized power & 0.292 [0.220, 0.370] & 0.259 [0.210, 0.323] & $1.1\times10^{-7}$ \\
Normalized jerk & 350.2 [283.3, 416.9] & 403.2 [303.3, 513.5] & 0.33 \\
\bottomrule
\end{tabular}
\end{table}

\paragraph{Result.} Table~\ref{tab:physical_analysis} shows no evidence for
the hypothesized interaction. For normalized torque and jerk, the
dynamic-class and quasi-static-class confidence intervals overlap and
Levene's test is not significant. Mass-normalized power shows a significant
Levene's $p$, but its clip-level bootstrap CIs overlap and the point
estimate is, if anything, marginally \emph{higher} for the dynamic classes
-- consistent with the row-level pseudo-replication described above rather
than a real, class-level effect (the same 128-replicates-per-clip
correlation that produced the misleading CV result). Small classes
(sitting $n{=}4$, kicking $n{=}7$, dancing $n{=}8$) further limit how finely
this test could distinguish individual classes, which is why dynamic and
quasi-static are pooled into two groups rather than compared class by
class.

We conclude that, for this checkpoint and this metric set, morphology's
effect on control demand does not detectably depend on the motion class
being performed: cross-shape variability is statistically indistinguishable
between highly dynamic and quasi-static motions once the near-zero-baseline
normalization artifact is corrected and inference is done at the clip
level. This is a properly negative result rather than an absence of
analysis: the instrumentation is validated (magnitudes match the
tracking-jerk figures already reported in this paper, and the effort-limit
normalization reuses the same mass-scaling formula validated elsewhere in
this work), and the same pipeline that would have reported a false positive
under a naive CV comparison instead surfaces its own artifact under a
correct, clip-level test. No mechanical or biomechanical claim is made
beyond what these instrumented rollouts show for this policy.
